\documentclass[10pt]{article}
\usepackage[margin=0.7in]{geometry}
\usepackage[T1]{fontenc}
\usepackage{lmodern}
\usepackage{microtype}
\usepackage{amsmath}
\usepackage{graphicx}
\usepackage{hyperref}
\hypersetup{colorlinks=true,urlcolor=blue,linkcolor=blue}
\setlength{\parindent}{0pt}
\setlength{\parskip}{0.35em}
\sloppy
\newcommand{\ananchor}[1]{\par{\footnotesize\emph{AN anchor: #1}}\par}

\title{\texorpdfstring{Barebones \(D^0\) Analysis Workflow}{Barebones D0 Analysis Workflow}}
\author{Evan Zhang}
\date{Codex export prepared 2026-06-18}

\begin{document}
\maketitle
\vspace{-1.4em}

\section*{Goal}

Reproduce the \(D^0\to K\pi\) mass peak starting from official MINIAOD. The
analysis note, official forests, existing skims, and final plots are validation
references, not the starting point for the reproduced chain.

\section*{Timeline}

\begin{enumerate}
\item \textbf{Freeze the MINIAOD input.}
Start from the 2023 HIForward MINIAOD datasets and the relevant certification
metadata. Save an explicit DAS file manifest for each CRAB production so every
forest can be traced back to exact MINIAOD files.
\ananchor{Sec. 2 lines 129--143 define the 2023 PbPb dataset, Golden JSON,
CMSSW/global-tag production context, run/luminosity table, trigger strategy,
and offline-selection references. Appendix A lines 996--1013 records the
run-range exclusions. The exact DAS manifest is reproduction bookkeeping added
here.}

\item \textbf{Forest the MINIAOD.}
Run CMSSW/CRAB foresting to reconstruct only the branches needed for this
analysis: event identifiers, vertex/filter information, tracks, PF/HF
candidates, ZDC rechits, and Dfinder \(K\pi\) candidates. For current large
tests, use raw CRAB output with minimal analyzer settings: keep analysis
content, disable irrelevant muon/PF-jet/ZDC-digi outputs where safe.
\ananchor{Sec. 2 lines 136--139 state the CMSSW/global-tag production context;
Sec. 7 lines 477--486 describes the \(D^0\) candidate construction from
oppositely charged tracks, the \(K/\pi\) mass hypotheses, and the secondary
vertex fit. The minimal branch pruning is our reproducibility implementation,
not a separate AN selection.}

\item \textbf{Validate the recreated forest.}
Before skimming, compare recreated forests against official forests/skims on
overlapping events where possible. Check schema, event counts, ZDC behavior,
and same-event Dfinder candidate behavior. If this fails, fix foresting before
touching the physics skim.
\ananchor{No standalone AN line: this is a reproduction closure check. The AN
targets being checked are the event selection in Sec. 5 lines 287--318, the
\(D^0\) reconstruction and candidate selection in Sec. 7 lines 477--504, and
the signal-extraction model in Sec. 7.2 lines 540--583.}

\item \textbf{Make the event skim.}
From the recreated forest, select good UPC-like events. Apply event-quality
filters, ZDC 0nXn topology selection, and HF rapidity-gap selection. Save both
the selected-event tree and all Dfinder candidates inside selected events.
\ananchor{Sec. 5.1 lines 287--318 define the offline event-selection strategy:
primary vertex, \(n_{\mathrm{Vtx}}\), halo rejection, ZDC XOR, and HF
rapidity-gap requirements. Sec. 5.2 lines 319--326 defines the offline ZDC
signal and 1n thresholds; Sec. 5.3 lines 349--366 motivates the HF \(E_{\max}\)
criterion, and lines 393--419 describe the scan and chosen HF thresholds.
Appendix E lines 1094--1124 gives the ZDC ADC-to-energy recipe.}

\emph{Cuts used:}
\begin{itemize}
\item certified lumisection selection from the Golden JSON;
\item trigger bookkeeping: \texttt{isZDCOr == 1};
\item primary-vertex boolean filter: \texttt{pprimaryVertexFilter == 1};
\item vertex-position quality cut: \(|z_{\mathrm{PV}}| < 15~\mathrm{cm}\);
\item vertex-multiplicity cut: \(1 \le \texttt{nVtx} \le 3\);
\item cluster-compatibility boolean filter:
\texttt{pclusterCompatibilityFilter == 1};
\item halo-rejection boolean filter: \texttt{cscTightHalo2015Filter == true};
\item derived ZDC cut:
\[
(\texttt{ZDCsumMinus} > T_{\mathrm{ZDC}}^- \ \&\&
 \texttt{ZDCsumPlus} < T_{\mathrm{ZDC}}^+)
\quad \mathrm{or} \quad
(\texttt{ZDCsumPlus} > T_{\mathrm{ZDC}}^+ \ \&\&
 \texttt{ZDCsumMinus} < T_{\mathrm{ZDC}}^-);
\]
\item derived HF cut:
\[
(\texttt{ZDCgammaN} \ \&\&\ \texttt{HFEMaxPlus} < T_{\mathrm{HF}}^+)
\quad \mathrm{or} \quad
(\texttt{ZDCNgamma} \ \&\&\ \texttt{HFEMaxMinus} < T_{\mathrm{HF}}^-).
\]
\end{itemize}

\item \textbf{Make the \(D^0\) candidate skim.}
From candidates in selected events, apply daughter-track quality cuts,
fiducial \(D^0\) kinematics, and bin-dependent topological cuts. Save the
selected \(D^0\) candidate tree.
\ananchor{Sec. 7 lines 477--504 defines \(D^0\) reconstruction, track criteria,
\(K/\pi\) assignment, secondary-vertex fitting, topological variables,
rectangular cuts, and the BDT cross-check. Lines 505--507 introduce the
before/after-BDT distribution checks; the MC/data topology comparisons are
shown around lines 573--585.}

\emph{Cuts used:}
\begin{itemize}
\item fit-range cut: \(1.70 < \texttt{Dmass} < 2.05~\mathrm{GeV}\);
\item \(D^0\) fiducial cuts:
\[
2 < \texttt{Dpt} \le 12~\mathrm{GeV},
\qquad |\texttt{Dy}| \le 2;
\]
\item daughter-track momentum cuts:
\[
\texttt{Dtrk1Pt} > 1~\mathrm{GeV},
\qquad \texttt{Dtrk2Pt} > 1~\mathrm{GeV};
\]
\item daughter-track acceptance cuts:
\[
|\texttt{Dtrk1Eta}| < 2.4,
\qquad |\texttt{Dtrk2Eta}| < 2.4;
\]
\item daughter-track relative momentum-error cuts:
\[
\texttt{Dtrk1PtErr}/\texttt{Dtrk1Pt} < 0.1,
\qquad
\texttt{Dtrk2PtErr}/\texttt{Dtrk2Pt} < 0.1;
\]
\item daughter-track high-purity boolean filters:
\texttt{Dtrk1highPurity == true} and \texttt{Dtrk2highPurity == true}, where
available in the from-forest branch set;
\item derived topology sideband cut used for optimization only:
\[
T_{\mathrm{sb}}^{\mathrm{in}}
< |\texttt{Dmass} - m_{D^0}|
< T_{\mathrm{sb}}^{\mathrm{out}};
\]
\item derived topological cuts:
These are the same four topological selections listed in AN Sec. 7
lines 487--493 and Table 7. The branch-to-AN mapping is:
\(\texttt{Dalpha}\) is the pointing angle \(\alpha\) between the reconstructed
\(D^0\) momentum and the primary-to-secondary vertex direction;
\(\texttt{DsvpvSig}\) is the normalized 3D decay length \(d_0/\sigma(d_0)\);
\(\texttt{Dchi2cl}\) is the secondary-vertex probability/confidence level
(not a raw \(\chi^2\)); and \(\texttt{Ddtheta}\) is the daughter-track opening
angle \(\Delta\theta\).
\[
\texttt{Dalpha} < T_{\alpha}(p_T,|y|), \qquad
\texttt{DsvpvSig} > T_{\mathrm{svpv}}(p_T,|y|),
\]
\[
\texttt{Dchi2cl} > T_{\mathrm{vtxProb}}(p_T,|y|), \qquad
\texttt{Ddtheta} < T_{\Delta\theta}(p_T,|y|).
\]
\end{itemize}

\item \textbf{Make the mass peak.}
Use the selected candidate tree to plot and fit the \(K\pi\) invariant-mass
distribution, inclusive and in the final analysis bins. The broad mass range is
for fitting and sidebands; it should not be tightened just to make a cleaner
peak.
\ananchor{Sec. 7.2 lines 540--583 defines raw-yield extraction, fit
components, MC-constrained swapped and peaking-background shapes, signal-width
handling, trigger choices by \(p_T\), and the binned \(D^0\) mass fits shown in
Figs. 34--36.}

\emph{Cuts used:}
\begin{itemize}
\item no additional peak-sculpting cut beyond the candidate skim;
\item final binning cuts in \(p_T\) and \(y\);
\item mass-fit model and sideband/cross-check definitions for yield extraction.
\end{itemize}
\end{enumerate}

\section*{Cut Families and How They Are Justified}

\begin{description}
\item[Event quality.]
Removes bad collision events, bad vertices, halo/noise-like events, and other
detector pathologies. These are standard reconstruction-quality requirements
from forest/skim branches, not cuts tuned on the \(D^0\) mass peak.

\item[ZDC topology.]
Selects UPC-like 0nXn events: one ZDC side is neutron-like and the other is
0n-like. The thresholds are data-driven. The method uses data/ZeroBias ZDC
energy structure and EmptyBX noise tails to identify the neutron onset while
controlling fake neutron assignment. The current official \(D^0\) MC is not
used for this because its checked detector-level ZDC neutron response is not
usable.

\item[HF rapidity gap.]
Rejects hadronic activity on the photon-going side. The threshold is derived
with an automated data/MC scan: prompt \(D^0\) MC for signal-like retention
where useful, low-track data as UPC-like control, high-track data as
hadronic-like control, and EmptyBX as a detector-noise cross-check.

\item[Daughter-track quality.]
Ensures the kaon/pion tracks are reconstructable and well measured. These cuts
use standard track quality, tracking acceptance, minimum track momentum, and
bounded relative momentum uncertainty. They are detector-quality cuts rather
than peak-optimized cuts.

\item[\(D^0\) fiducial region.]
Defines the analysis phase space in \(p_T\) and \(y\). This is part of the
measurement definition. The broad mass range is used for fitting and sideband
studies.

\item[Topological cuts.]
Suppress random \(K\pi\) combinations while retaining displaced,
vertex-compatible \(D^0\) decays. The method uses prompt \(D^0\) MC for signal
efficiency and data sidebands for combinatorial background. Current production
uses rectangular cuts in pointing angle, decay-length significance, vertex
probability, and angular consistency, chosen by a predeclared
efficiency/background plateau procedure.
\end{description}

\section*{How the Nontrivial Cuts Are Generated}

\begin{description}
\item[ZDC thresholds.]
Characterize the ZDC neutron-energy structure in data and compare it with
EmptyBX/noise tails. Choose a stable neutron-onset threshold with controlled
noise leakage. This is explicitly data-driven.

\item[HF thresholds.]
Scan the gap-side maximum HF energy. Require high retention for signal-like or
UPC-like samples and rejection of high-activity data. EmptyBX is used to check
that the chosen threshold is not a detector-noise artifact.

\item[Topology sidebands.]
Use MC mass resolution and data sideband balance to define sidebands away from
the \(D^0\) signal peak. This prevents optimizing the cuts directly on the
signal region.

\item[Topology floors.]
First derive basic displaced-vertex quality floors with loose pointing cuts.
Use MC signal efficiency plus data-sideband and nonmatched-MC background
rejection.

\item[Pointing cuts.]
With the floors fixed, scan pointing variables and accept a stable plateau of
choices that preserves signal and does not increase background proxies. The
analysis-note row is treated as a closure point inside the plateau, not as the
hidden optimizer.
\end{description}

\section*{Example Cut Diagrams}

\textbf{ZDC thresholds.}

\begin{center}
\begin{minipage}{0.48\textwidth}
\centering
\includegraphics[width=\textwidth]{support/figures/preflight_medium_scale/zdc_minus_1n_threshold.png}
\end{minipage}
\hfill
\begin{minipage}{0.48\textwidth}
\centering
\includegraphics[width=\textwidth]{support/figures/preflight_medium_scale/zdc_plus_1n_threshold.png}
\end{minipage}
\end{center}

\emph{How the cut is made.}
The ZDC energy spectrum is compared with EmptyBX/noise behavior. A neutron
threshold is placed at the stable onset of the one-neutron structure, far
enough above the noise tail to keep fake neutron assignments controlled. Events
pass the 0nXn topology when one side is neutron-like and the opposite side is
consistent with 0n/noise.

\textbf{HF thresholds.}

\begin{center}
\begin{minipage}{0.48\textwidth}
\centering
\includegraphics[width=\textwidth]{support/figures/hf_gap_derivation/hf_minus_threshold_efficiencies.png}
\end{minipage}
\hfill
\begin{minipage}{0.48\textwidth}
\centering
\includegraphics[width=\textwidth]{support/figures/hf_gap_derivation/hf_plus_threshold_efficiencies.png}
\end{minipage}
\end{center}

\emph{How the cut is made.}
The scanned variable is the maximum HF energy on the side that is supposed to
be quiet. For a proposed threshold \(T\), the event passes the HF rapidity-gap
cut if
\[
E_{\mathrm{HF,max}}^{\mathrm{gap}} < T.
\]
The \(nTrack\) requirement is not applied as a physics selection. It is used
only to build data control samples for choosing \(T\). The low/high
\(\mathrm{nTrack}\) samples are built from
\(\mathrm{nTrackInAcceptanceHP}\): low \(\mathrm{nTrack}\) means
\(\mathrm{nTrack} \le 3\), and high \(\mathrm{nTrack}\) means
\(\mathrm{nTrack} \ge 10\), for the current full-sample derivation. The
low-\(\mathrm{nTrack}\) pass curve is
\[
\mathrm{Pass}_{\mathrm{low}}(T)
=
\frac{N(\mathrm{nTrack} \le 3,\ E_{\mathrm{HF,max}}^{\mathrm{gap}}<T)}
{N(\mathrm{nTrack} \le 3)},
\]
which is a UPC-like data-retention proxy. The high-\(\mathrm{nTrack}\) leak
curve is
\[
\mathrm{Leak}_{\mathrm{high}}(T)
=
\frac{N(\mathrm{nTrack} \ge 10,\ E_{\mathrm{HF,max}}^{\mathrm{gap}}<T)}
{N(\mathrm{nTrack} \ge 10)},
\]
which is a hadronic-activity leakage proxy. The current working numerical rule
is
\[
\mathrm{Pass}_{\mathrm{low}}(T) \ge 0.9765,
\qquad
\mathrm{Leak}_{\mathrm{high}}(T) \le 0.85,
\]
and, when a same-side prompt \(D^0\) MC control sample is available,
\[
\epsilon_{\mathrm{MC}}^{\mathrm{gap}}(T) \ge 0.989.
\]
For each side, the operational choice is the smallest threshold satisfying
these criteria on a stable plateau, with EmptyBX used to check that the choice
is not dominated by detector noise. When the MC control sample is absent, the
MC efficiency is recorded as unavailable rather than as a physical zero. In the
plots, the selected point is therefore not the value that makes the prettiest
mass peak; it is the value where the retention/rejection curves show a stable
rapidity-gap definition.

\textbf{Topology sidebands.}

\begin{center}
\includegraphics[width=0.72\textwidth]{support/figures/topology_sideband_independent/sideband_resolution_rule.png}
\end{center}

\emph{How the cut is made.}
The mass sideband is one global window, not a different window for each
\(p_T,|y|\) bin:
\[
0.07 < |\texttt{Dmass} - 1.86484| < 0.12~\mathrm{GeV}.
\]
Equivalently, the lower and upper sidebands are
\[
1.74484 < \texttt{Dmass} < 1.79484,
\qquad
1.93484 < \texttt{Dmass} < 1.98484.
\]
This same mass window is applied separately inside each analysis bin. The
topological cut values are bin-dependent, but the sideband mass definition is
global. The window is derived from prompt truth-matched MC: estimate the
mass-resolution scale \(\sigma_{68}\) in each bin, take the worst bin, and set
the inner and outer sideband edges near \(3\sigma_{\max}\) and
\(5\sigma_{\max}\), rounded to 10 MeV. The plotted points are therefore the MC
mass-resolution scales in each bin; the horizontal lines are the resulting
global sideband edges. These sideband candidates act as the data-background
proxy for topological scans, so the signal region is not used to choose the
cuts.

\textbf{Topology floors.}

\begin{center}
\emph{\(\texttt{Dchi2cl}\) floor scan.}\\[0.2em]
\includegraphics[width=0.86\textwidth]{support/figures/topological_floor_independent_sideband/chi2_floor_profile_grid.png}

\vspace{0.5em}
\emph{\(\texttt{DsvpvSig}\) floor scan after fixing \(\texttt{Dchi2cl}\).}\\[0.2em]
\includegraphics[width=0.86\textwidth]{support/figures/topological_floor_independent_sideband/floor_profile_grid.png}
\end{center}

\emph{How the cut is made.}
The floors are lower bounds on the basic displaced-vertex quality variables:
\[
\texttt{DsvpvSig} > T_{\mathrm{svpv}},
\qquad
\texttt{Dchi2cl} > T_{\mathrm{vtxProb}}.
\]
\(\texttt{DsvpvSig}\) is the secondary-to-primary vertex displacement divided
by its uncertainty, so larger values mean a more clearly displaced \(D^0\)
decay. \(\texttt{Dchi2cl}\) is the secondary-vertex fit confidence, so larger
values mean the two tracks form a more credible common decay vertex. This is
the AN's ``Secondary vertex probability'' cut in Table 7, despite the branch
name containing ``chi2cl''.

The scan uses loose pointing cuts so the floor variables are isolated. Prompt
truth-matched MC is the signal proxy; non-truth-matched MC and data sidebands
are background proxies. The current rule first chooses the first
\(\texttt{Dchi2cl}\) floor with MC-background rejection \(\ge 0.04\) and MC
signal loss \(\le 0.06\). It then chooses the first \(\texttt{DsvpvSig}\)
floor with MC-background rejection \(\ge 1/3\), data-sideband rejection
\(\ge 0.18\), and MC signal retention \(\ge 0.55\). This gives
\(\texttt{Dchi2cl}>0.1\) in all bins, \(\texttt{DsvpvSig}>2.5\) for
\(2<p_T<5\), and \(\texttt{DsvpvSig}>3.5\) for higher \(p_T\). In the plots,
the x-axis is the scanned floor variable, the curves show signal retention and
background rejection, and the dashed line is the selected floor.

\textbf{Pointing cuts.}

\begin{center}
\includegraphics[width=0.86\textwidth]{support/figures/topological_pointing_independent_sideband/pointing_plateau_grid.png}
\end{center}

\emph{How the cut is made.}
After the topology floors are fixed, the pointing variables are scanned as
upper bounds: smaller angles mean the reconstructed \(D^0\) momentum points
more consistently back to the production vertex. The selected value is taken
from a plateau where MC signal efficiency remains high and the sideband
background proxy does not rise. This separates the pointing choice from the
floor choice and avoids tuning all topology variables directly on the signal
peak. In AN language, this section covers the pointing angle \(\alpha\) and
the daughter-track opening angle \(\Delta\theta\).

\section*{BDT Topology Optimization Check}

A ROOT/TMVA BDT was trained as a cross-check using the same topology variables:
\[
\texttt{Dalpha},\quad \texttt{DsvpvSig},\quad \texttt{Dchi2cl},\quad
\texttt{Ddtheta}.
\]
These are the AN variables \(\alpha\), normalized 3D decay length,
secondary-vertex probability, and \(\Delta\theta\), respectively.
Signal was prompt truth-matched \(D^0\) MC. Background was data sidebands.
Two threshold rules were tested: a peak-blind Punzi scan on held-out MC and
sidebands, and a data-split scan where half the data chooses the threshold by
sideband-subtracted peak significance and the other half validates it.

\begin{center}
\small
\begin{tabular}{lrrrrr}
\hline
Bin & BDT cut & MC eff. & Sideband rej. & Val. \(Z\) & AN \(Z\) \\
\hline
\(|y|<1,\ 2<p_T<5\) & \(-0.950\) & 1.000 & 0.000 & 2.67 & 4.61 \\
\(|y|<1,\ 5<p_T<8\) & \(-0.950\) & 1.000 & 0.000 & 5.16 & 6.33 \\
\(|y|<1,\ 8<p_T<12\) & \(-0.950\) & 1.000 & 0.000 & 2.68 & 3.08 \\
\(1<|y|<2,\ 2<p_T<5\) & \(-0.950\) & 1.000 & 0.000 & 0.26 & 2.71 \\
\(1<|y|<2,\ 5<p_T<8\) & \(-0.950\) & 1.000 & 0.000 & 2.26 & 2.86 \\
\(1<|y|<2,\ 8<p_T<12\) & \(-0.950\) & 1.000 & 0.000 & 1.40 & 4.73 \\
\hline
\end{tabular}
\end{center}

The left edge of the scan is \(-0.950\), so the selected BDT threshold accepts
all candidates in every bin. In this protocol, the BDT does not find a
nontrivial yield-optimizing topology cut. The nominal rectangular topology
selection remains the more useful working choice for producing a clean peak.

\begin{center}
\begin{minipage}{0.48\textwidth}
\centering
\includegraphics[width=\textwidth]{support/figures/bdt_topology_98forests/absy0to1_pt2to5_bdt_threshold_scan.png}
\end{minipage}
\hfill
\begin{minipage}{0.48\textwidth}
\centering
\includegraphics[width=\textwidth]{support/figures/bdt_topology_98forests/absy0to1_pt5to8_bdt_threshold_scan.png}
\end{minipage}

\vspace{0.4em}
\begin{minipage}{0.48\textwidth}
\centering
\includegraphics[width=\textwidth]{support/figures/bdt_topology_98forests/absy0to1_pt8to12_bdt_threshold_scan.png}
\end{minipage}
\hfill
\begin{minipage}{0.48\textwidth}
\centering
\includegraphics[width=\textwidth]{support/figures/bdt_topology_98forests/absy1to2_pt2to5_bdt_threshold_scan.png}
\end{minipage}

\vspace{0.4em}
\begin{minipage}{0.48\textwidth}
\centering
\includegraphics[width=\textwidth]{support/figures/bdt_topology_98forests/absy1to2_pt5to8_bdt_threshold_scan.png}
\end{minipage}
\hfill
\begin{minipage}{0.48\textwidth}
\centering
\includegraphics[width=\textwidth]{support/figures/bdt_topology_98forests/absy1to2_pt8to12_bdt_threshold_scan.png}
\end{minipage}
\end{center}

\begin{center}
\begin{minipage}{0.48\textwidth}
\centering
\includegraphics[width=\textwidth]{support/figures/bdt_topology_98forests/absy0to1_pt5to8_mass_bdt_vs_an.png}
\end{minipage}
\hfill
\begin{minipage}{0.48\textwidth}
\centering
\includegraphics[width=\textwidth]{support/figures/bdt_topology_98forests/absy1to2_pt8to12_mass_bdt_vs_an.png}
\end{minipage}
\end{center}

\section*{Extra Reproduction Work Beyond the Analysis Note}

\begin{itemize}
\item Exact CRAB input manifests are saved for every production.
\item Forest output is minimized to the branch set needed for this \(D^0\)
workflow where safe.
\item Recreated forests are validated against official products on same events
before large-scale use.
\item ZDC reconstruction and calibration/backport behavior are cross-checked
against official reforested outputs.
\item ZDC, HF, sideband, and topology choices are rederived or independently
justified using data, control samples, MC, and predeclared criteria.
\item CRAB smoke and medium-scale tests are used before full production to
check job splitting, event counts, output transfer, schema, and skim-to-peak
plumbing.
\item Task names, paths, input manifests, and validation results are logged so
the chain is auditable.
\end{itemize}

\section*{One-Line Chain}

\[
\text{MINIAOD}
\rightarrow \text{recreated forest}
\rightarrow \text{event skim}
\rightarrow \text{\(D^0\) candidate skim}
\rightarrow \text{\(K\pi\) mass peak}.
\]

\end{document}
